\chapter{Metodología}
\label{cap:metodologia}

Este capítulo describe la metodología seguida para el desarrollo del
Trabajo Fin de Grado. El trabajo tiene dos caras. Por un lado, es un
proyecto de desarrollo software, en el que se construye una herramienta
a partir de unos requisitos; por otro, incluye una parte experimental
orientada a comprobar si el sistema resuelve el problema planteado. Esta
dualidad condiciona el enfoque metodológico, que se detalla
en la sección~\ref{sec:enfoque-metodologico}, y se concreta en las fases
del proyecto (sección~\ref{sec:fases-proyecto}), la metodología de
desarrollo (sección~\ref{sec:metodologia-desarrollo}), las herramientas
empleadas (sección~\ref{sec:herramientas-tecnologias}) y la metodología
de evaluación (sección~\ref{sec:metodologia-evaluacion}).


\section{Enfoque metodológico}
\label{sec:enfoque-metodologico}

El desarrollo de la herramienta se ha abordado siguiendo un modelo
iterativo e incremental. En lugar de planificar el sistema completo de
antemano y construirlo en una única fase, se ha optado por levantar un
prototipo mínimo funcional y ampliarlo de forma progresiva, validando su
funcionamiento de extremo a extremo en cada iteración antes de añadir
nueva complejidad. Este enfoque resulta especialmente adecuado ante la incertidumbre
técnica que supone el uso de modelos de lenguaje de gran tamaño, cuyo
comportamiento no puede preverse por completo sin experimentación.

Al carácter de desarrollo software se suma un componente experimental.
Una vez construido el sistema, no basta con comprobar que funciona desde
un punto de vista técnico: es necesario valorar si las preguntas que
genera son realmente útiles para un docente y si las decisiones de
diseño adoptadas están justificadas. Por este motivo, la metodología
contempla además una fase de evaluación con la que contrastar
empíricamente la aportación de los distintos componentes del sistema.


\section{Fases del proyecto}
\label{sec:fases-proyecto}

El trabajo se ha organizado en cuatro fases principales que, si bien se
describen de forma secuencial, se han solapado parcialmente en el
tiempo, en coherencia con el carácter iterativo del desarrollo:

\begin{enumerate}
  \item \textbf{Revisión bibliográfica.} Estudio del estado del arte en
        generación automática de preguntas con modelos de lenguaje, en
        evaluación auténtica frente a la aparición de la inteligencia
        artificial generativa y en los antecedentes más próximos al
        problema abordado. Esta fase, recogida en el
        capítulo~\ref{cap:estado-arte}, ha permitido delimitar la
        aportación del trabajo.
  \item \textbf{Diseño de la arquitectura.} Definición de los
        componentes del sistema y de sus interfaces: el módulo de
        análisis de código, el sistema multi-agente de generación y
        revisión, y la configuración académica que adapta las preguntas.
  \item \textbf{Implementación iterativa.} Construcción incremental de
        los componentes, partiendo de un prototipo capaz de generar una
        pregunta sencilla y ampliándolo de forma progresiva con nuevas
        capacidades, validando cada incremento de extremo a extremo.
  \item \textbf{Evaluación.} Diseño y ejecución de las pruebas
        destinadas a valorar la calidad de las preguntas generadas y la
        aportación del agente revisor, según la metodología descrita en
        la sección~\ref{sec:metodologia-evaluacion}.
\end{enumerate}


\section{Metodología de desarrollo}
\label{sec:metodologia-desarrollo}

La implementación se ha guiado por el principio de mantener en todo
momento una versión del sistema ejecutable y verificable. Cada
incremento ha consistido en una funcionalidad concreta y acotada ---por
ejemplo, la extracción de unidades de código, el filtrado de las
triviales o la incorporación de un nuevo tipo de pregunta--- que se ha
integrado y verificado de extremo a extremo antes de continuar. Esta forma de trabajar permite
detectar errores de forma temprana y reduce el riesgo asociado a la
integración de componentes que dependen de un modelo de lenguaje
externo.

El control de versiones se ha gestionado con Git. El proyecto se ha
estructurado en dos repositorios independientes: uno para el código de
la herramienta y otro para la presente memoria. Esta separación facilita
el seguimiento de la evolución de cada parte de forma autónoma y
mantiene un historial de cambios claro y diferenciado.


\section{Herramientas y tecnologías}
\label{sec:herramientas-tecnologias}

El desarrollo se ha llevado a cabo en el lenguaje Python, por su amplio
ecosistema en el ámbito del procesamiento del lenguaje natural y por la
disponibilidad de las bibliotecas necesarias para el proyecto. El
análisis sintáctico del código fuente se apoya en \emph{tree-sitter},
una biblioteca de análisis incremental que ofrece gramáticas para
múltiples lenguajes de programación y que dota al sistema de capacidad
de extensión a lenguajes distintos del empleado en la validación
inicial.

La orquestación del sistema multi-agente se ha realizado con el
\emph{framework} CrewAI, que permite definir agentes con funciones
diferenciadas y coordinar su interacción. El modelo de lenguaje se
consume a través de Amazon Bedrock, empleando Claude durante el
desarrollo en su variante más ligera por motivos de coste y rapidez, y
reservando un modelo de mayor capacidad para la validación final. La
redacción de la memoria se ha realizado en \LaTeX. La
tabla~\ref{tab:herramientas} resume las principales herramientas
empleadas y su finalidad.

\begin{table}[htbp]
  \centering
  \caption{Principales herramientas y tecnologías empleadas.}
  \label{tab:herramientas}
  \begin{tabular}{ll}
    \toprule
    Herramienta & Finalidad \\
    \midrule
    Python        & Lenguaje de desarrollo de la herramienta \\
    tree-sitter   & Análisis sintáctico del código fuente \\
    CrewAI        & Orquestación del sistema multi-agente \\
    Amazon Bedrock & Acceso a los modelos de lenguaje \\
    Git           & Control de versiones \\
    \LaTeX        & Redacción de la memoria \\
    \bottomrule
  \end{tabular}
\end{table}


\section{Metodología de evaluación}
\label{sec:metodologia-evaluacion}

La evaluación del sistema se plantea en torno a una validación central,
alineada con el objetivo del trabajo, y un estudio de apoyo que justifica
una de sus decisiones de diseño. La validación central
(sección~\ref{sec:validacion-utilidad}) busca determinar si la
herramienta cumple su propósito: generar preguntas útiles para el
docente. El estudio de apoyo (sección~\ref{sec:ablacion}) es un estudio
de ablación que mide la aportación concreta del agente revisor.

\subsection{Validación de utilidad y configurabilidad}
\label{sec:validacion-utilidad}

La validación central consiste en valorar si las preguntas generadas
serían realmente utilizables por un docente en un examen, lo que
constituye la medida última del valor de la herramienta frente al
problema que la motiva. Para ello se analiza el resultado que produce el
sistema sobre distintos repositorios de código, examinando la
corrección técnica de las preguntas, su relevancia respecto al código
del que parten y su adecuación al nivel académico indicado.

Esta valoración se apoya en el criterio del tutor del trabajo, docente
en activo, cuya revisión de una muestra de las preguntas generadas
aporta un juicio de utilidad fundamentado desde la experiencia real en
el aula. Asimismo, se ilustra de forma cualitativa la
configurabilidad del sistema, mostrando cómo la variación de parámetros
como el nivel académico, el tipo de pregunta o el contexto de la
asignatura repercute en las preguntas producidas.

\subsection{Estudio de ablación del agente revisor}
\label{sec:ablacion}

Como estudio complementario se plantea una ablación que justifica la
decisión de incorporar un agente revisor al sistema. El experimento
compara el sistema completo con una variante en la que se desactiva
dicho agente, dejando únicamente la generación directa de preguntas, con
el fin de determinar en qué medida la revisión automática mejora la
calidad del resultado respecto a una generación sin postprocesado. Con
este fin, se generan preguntas sobre un mismo conjunto de código en
ambas configuraciones y se comparan los resultados. La comparación se
realiza con una rúbrica que valora aspectos como la corrección técnica,
la claridad del enunciado y la adecuación al nivel. Para aligerar la
evaluación manual, se contempla emplear además un modelo de lenguaje
externo como juez automático. Conviene aclarar que este juez es
independiente del agente revisor: el revisor forma parte del proceso de
generación de preguntas, mientras que el juez interviene solo en la fase
de evaluación, como evaluador imparcial de los resultados.